\section{Source Section III.C: Piston-Effect Scenario}
\label{sec:piston-effect}

\sourceeqrange{Eqs. (3.13)--(3.16) and figures 3--4}

\subsection*{Purpose}

This subsection tracks the analytical estimates used to read the bubble-in-liquid
piston-effect example.  It separates pressure and temperature estimates from the
computed flow field.

\subsection*{Reading Guide}

The scenario uses the bottom/top temperature branch under zero gravity.  The
bubble radius, wall temperature change, pressure increment, and adiabatic
coefficient are source inputs for Eqs. (3.13)--(3.15).  Equation (3.16) then
introduces a latent-heat velocity scale, which is useful for reading interface
mass transfer and flow magnitudes.  These formulas organize interpretation; they
are not a reduced simulation model.

\subsection*{Finite Estimate}

The checkable scalar relations are the pressure estimate proportional to the
wall-temperature change, the phase-temperature difference proportional to the
phase difference in adiabatic coefficients, and the units of the latent-heat
velocity \(v_c=Q/(n_gT\Delta s)\).  The companion does not calculate the
reported time traces or velocity field.

The adiabatic coefficient is defined by
\begin{equation}
  A_s=\left(\frac{\partial T}{\partial p}\right)_s.
  \sourceeqbadge{3.14}
\end{equation}
At fixed entropy per particle, the Section II.A thermodynamic derivatives give
\begin{equation}
  \left(\frac{\partial T}{\partial n}\right)_s
  =\frac{2T}{dn(1-v_0n)},
  \qquad
  \left(\frac{\partial p}{\partial n}\right)_s
  =\frac{k_BT}{(1-v_0n)^2}
    \left(1+\frac2d-\frac{T_s}{T}\right).
\end{equation}
Thus
\begin{align}
  A_s
  &=\frac{(\partial T/\partial n)_s}{(\partial p/\partial n)_s} \\
  &=\frac{2(1-v_0n)}{d n k_B(1+2/d-T_s/T)}.
\end{align}
For the two-dimensional simulations, \(d=2\), so
\begin{equation}
  A_s=\frac{1/n-v_0}{k_B(2-T_s/T)}.
  \sourceeqbadge{3.14}
  \label{eq:iii-piston-As-derived}
\end{equation}
This connects the piston-effect estimate back to the same Section II.A
thermodynamic branch used for Eq. (3.9).

For a small adiabatic pressure perturbation, \(\delta T\approx A_s\delta p\).
Using the liquid value for the surrounding phase gives the source pressure
scale
\begin{equation}
  \delta p\sim A_{s\ell}^{-1}\delta T_b.
  \sourceeqbadge{3.13}
\end{equation}
If a common pressure perturbation is read in the two phases, then
\begin{equation}
  \delta T_g=A_{sg}\delta p,
  \qquad
  \delta T_\ell=A_{s\ell}\delta p,
\end{equation}
and subtraction gives
\begin{equation}
  (\Delta T)_{g\ell}=(A_{sg}-A_{s\ell})\delta p.
  \sourceeqbadge{3.15}
\end{equation}

Finally, Eq. (3.16) balances heat flux against latent heat per unit volume.  If
\(\Delta s=s_g-s_\ell\), the latent heat per particle is \(T\Delta s\), and a
normal particle flux \(n_gv_c\) carries heat flux \(n_gv_cT\Delta s\).  Solving
gives
\begin{equation}
  v_c=\frac{Q}{n_gT\Delta s}
  =-\frac{\lambda_\ell T'_\infty}{n_gT\Delta s}.
  \sourceeqbadge{3.16}
  \label{eq:iii-latent-heat-velocity-derived}
\end{equation}
The second equality is the far-field heat-flux branch
\(Q=-\lambda_\ell T'_\infty\).
Here \(Q\) is an energy flux, while \(n_gT\Delta s\) has energy per volume;
therefore the quotient has units of velocity.  This confirms the scale used to
read the source figures, not the flow field itself.
The finite pressure/temperature-estimate and latent-velocity checks are
covered by the Section III mirror pair listed in
Section~\ref{sec:numerical-results-overview}; they do not compute the source
velocity field.

\subsection*{Limitations}

The local/global energy map, boundary heat transfer, and wall surface terms stay
external to this scenario guide.  No conclusion about full continuum-model
equivalence follows from these scalar estimates.
